\section{Conjectures and Theorems}

\subsection{Basic Definitions}

\begin{definition}
A \emph{conjecture} is a mathematical statement that is believed to be true
based on evidence or intuition, but which has not yet been proved or disproved.
\end{definition}

\begin{definition}
A \emph{theorem} is a mathematical statement that has been rigorously proved
to be true from previously established results (such as axioms, definitions,
and other theorems).
\end{definition}

\subsection{Examples from Number Theory}

\begin{example}[A true conjecture that became a theorem]
Historically, mathematicians \emph{conjectured} that there are infinitely many prime numbers.
This was later proved by Euclid and is now a theorem:

\begin{theorem}[Euclid]
There are infinitely many prime numbers.
\end{theorem}

Thus, a statement can start as a conjecture and, once proved, become a theorem.
\end{example}

\begin{example}[A conjecture that is still open]
\emph{Goldbach's Conjecture} is one of the most famous unsolved problems in number theory:

\begin{conjecture}[Goldbach]
Every even integer \( n \geq 4 \) can be written as the sum of two prime numbers.
\end{conjecture}

This statement has been checked for very large values of \( n \), but no complete proof
or counterexample is known.
\end{example}

\begin{example}[A false conjecture]
Consider the following statement:

\begin{conjecture}
Every odd integer \( n > 1 \) is a prime number.
\end{conjecture}

At first glance, small examples like \( 3, 5, 7 \) might suggest this is true.
However, it is \emph{false}. For instance,


\[
9 = 3 \cdot 3, \quad 15 = 3 \cdot 5,
\]


so \( 9 \) and \( 15 \) are odd but not prime. These are counterexamples that disprove
the conjecture.
\end{example}

\subsection{Remarks}

\begin{remark}
Conjectures play a central role in mathematical research. They guide exploration and
often lead to new techniques, whether they eventually become theorems or are disproved
by counterexamples.
\end{remark}
